%! Populations, Samples, Parameters, and Statistics — Unit 1, Lesson 1.2 — Group Activity (Answer Key)
\documentclass[10pt]{article}
\usepackage{saar-article}
\usepackage{saar-key}   % pulls in -boxes; do NOT also load -boxes
\newcommand{\TallMath}[1]{$\displaystyle #1\rule[-1.4em]{0pt}{3.2em}$}

\begin{document}

\pageheader{Unit 1, Lesson 1.2}{Group Activity --- Answer Key}
% NAMESTRIP: no \namedateperiod here — mirrors the blank. Teacher-only prose goes
% in the lesson plan as \begin{teachernote}[Group Activity], never in a key.

\vspace{0.15in}

\begin{headlinebox}{frost}
\small
\textbf{Who Are You Talking About? --- The Millbrook Public Library.} The library
has $1{,}200$ card holders and one question: what do people mainly use it for?
Every answer needs a \emph{reason}, and every number has to be labeled ---
parameter or statistic.
\end{headlinebox}

\vspace{0.1in}

\boxguard[30]
\begin{tcolorbox}[colback=white, colframe=black!40, breakable,
    title=\textbf{Tier R --- Remediate},
    fonttitle=\bfseries, arc=2mm, left=3mm, right=3mm, top=2mm, bottom=2mm]
\small
On a Saturday, a librarian asked $60$ card holders at the front desk, ``What do
you mainly use the library for?''

\begin{enumerate}[leftmargin=1.6em, itemsep=6pt, topsep=4pt]

  \item Name the two groups in this study.

        Population: \ans{all 1,200 card holders} \hfill Sample size: \ans{60}

        Sample: \ans{the 60 card holders asked at the desk}

  \item Complete the frequency table. The first row is done for you.

        \begin{center}
        \renewcommand{\arraystretch}{1.4}
        \begin{tabularx}{0.9\linewidth}{@{} l c X @{}}
        \toprule
        \textbf{Main use} & \textbf{Count} &
        \textbf{Relative frequency (percent)} \\
        \midrule
        Borrowing books & 21 & $21/60 = 0.35 = 35\%$ \\
        Using computers & 18 & \ans{$18/60 = 0.30 = 30\%$} \\
        Attending programs & 15 & \ans{$15/60 = 0.25 = 25\%$} \\
        Using study rooms &  6 & \ans{$6/60 = 0.10 = 10\%$} \\
        \bottomrule
        \end{tabularx}
        \end{center}

  \item Label each number \textbf{parameter} or \textbf{statistic}.

        \begin{center}
        \renewcommand{\arraystretch}{1.4}
        \begin{tabularx}{0.95\linewidth}{@{} X c @{}}
        \toprule
        \textbf{The number} & \textbf{Which kind?} \\
        \midrule
        The library has $1{,}200$ card holders.                    & \ans{parameter} \\
        $35\%$ of the $60$ card holders asked mainly borrow books. & \ans{statistic} \\
        \bottomrule
        \end{tabularx}
        \end{center}

  \item Estimate how many of the $1{,}200$ card holders mainly borrow books.

        \begin{work}
          0.35 \times 1200 &= 420 \text{ card holders}
        \end{work}

  \item Write one sentence reporting what the survey found. Name the group and
        the percent.
        \ansline{Of the 60 card holders asked, 35\% said they mainly borrow books.}
\end{enumerate}
\end{tcolorbox}

\vspace{0.1in}

\boxguard
\begin{tcolorbox}[colback=white, colframe=black!40, breakable,
    title=\textbf{Tier A --- Approaching Proficiency},
    fonttitle=\bfseries, arc=2mm, left=3mm, right=3mm, top=2mm, bottom=2mm]
\small
\begin{enumerate}[leftmargin=1.6em, itemsep=5pt, topsep=3pt]

  \item The librarian did not survey all $1{,}200$ card holders. Name the
        constraint in each situation.

        \begin{center}
        \renewcommand{\arraystretch}{1.45}
        \begin{tabularx}{0.97\linewidth}{@{} X p{3cm} @{}}
        \toprule
        \textbf{Situation} & \textbf{Constraint} \\
        \midrule
        Many card holders have moved and the library has no working phone number
        for them. & \ans{access} \\
        Reaching everyone by mail would cost more than the library's whole survey
        budget. & \ans{cost} \\
        New people get cards and others let theirs lapse every week.
          & \ans{change} \\
        \bottomrule
        \end{tabularx}
        \end{center}

  \item The next Saturday a second librarian asked a \emph{different} $60$ card
        holders the same question. This time $27$ said they mainly borrow books.

        \textbf{(a)} What percent of that second sample is that?

        \begin{work}
          \dfrac{27}{60} &= 0.45 \\
                         &= 45\%
        \end{work}

        \textbf{(b)} The two surveys gave $35\%$ and $45\%$. Did one of the
        librarians make a mistake? Explain.
        \ansline{No. Different samples give different statistics --- that is sampling variability.}

        \textbf{(c)} How many true percents are there for all $1{,}200$ card
        holders? \ans{1}

  \item Six card holders reported how many items they checked out on their last
        visit: $4, \; 7, \; 2, \; 5, \; 9, \; 3$.

        \textbf{(a)} Find the mean number of items.

        \begin{work}
          \bar{x} &= \dfrac{4+7+2+5+9+3}{6} \\
                  &= \dfrac{30}{6} \\
                  &= 5 \text{ items}
        \end{work}

        \textbf{(b)} Is that $5$ a parameter or a statistic? Say how you know.
        \ansline{A statistic --- it describes the six card holders asked, not all 1,200.}
\end{enumerate}
\end{tcolorbox}

\vspace{0.1in}

\boxguard
\begin{tcolorbox}[colback=white, colframe=black!40, breakable,
    title=\textbf{Tier E --- Extension},
    fonttitle=\bfseries, arc=2mm, left=3mm, right=3mm, top=2mm, bottom=2mm]
\small
\begin{enumerate}[leftmargin=1.6em, itemsep=5pt, topsep=3pt]

  \item The library director drafts this line for the annual report:

        \begin{center}
        \textit{``Millbrook's $1{,}200$ card holders mainly borrow books.''}
        \end{center}

        The $35\%$ behind it is computed correctly. Explain \emph{two} separate
        things this sentence claims that the Tier R data cannot support.
        \ansline{The 35\% describes only the 60 people asked, not all 1,200 card holders.}
        \ansline{And 35\% is not most of them --- 65\% of that sample said something else.}

  \item Millbrook has $1{,}200$ card holders; the Ashford branch across town has
        $400$. Each library surveys $60$ of its own card holders, and each gets
        $35\%$ for ``mainly borrows books.''

        Do those two $35\%$ figures estimate the \emph{same} parameter? Explain.
        \ansline{No. The two samples come from two different populations, so each 35\%}
        \ansline{estimates its own branch's true percent. Equal statistics, different parameters.}

  \item Your group is asked to estimate the mean number of minutes a Millbrook
        card holder spends in the library per visit. Plan it.

        \begin{center}
        \renewcommand{\arraystretch}{1.35}
        \begin{tabularx}{0.97\linewidth}{@{} p{4.2cm} X @{}}
        \toprule
        Our population & \ans{All 1,200 Millbrook card holders.} \\
        Our sample, and its size & \ans{Answers vary --- a stated group and a stated number.} \\
        The statistic we will compute & \ans{The mean minutes per visit for the people we timed.} \\
        The parameter it estimates & \ans{The mean minutes per visit for all 1,200 card holders.} \\
        One constraint we face & \ans{Answers vary --- time, cost, access, or change.} \\
        \bottomrule
        \end{tabularx}
        \end{center}
\end{enumerate}
\end{tcolorbox}

\end{document}
